\documentclass[12pt,a4paper]{article}

% ==========================================
% 1. ENCODAGE ET POLICES (FONTS)
% ==========================================
\usepackage[T1]{fontenc}   % Encodage des polices de sortie pour la gestion des accents
\usepackage{times}        % Police Times New Roman pour le texte
\usepackage{lmodern}      % Police Latin Modern améliorée
\usepackage{pifont}       % Symboles et caractères spéciaux (ex: \ding{45})

% ==========================================
% 2. DISPOSITION DE LA PAGE ET EN-TÊTES/PIEDS DE PAGE
% ==========================================
\usepackage[left=1cm, right=1cm, top=1cm, bottom=1.7cm]{geometry} % Configuration des marges
\usepackage{fancyhdr}     % Personnalisation des en-têtes et pieds de page
\usepackage{lastpage}     % Accès au numéro de la dernière page pour "page X/Y"
\usepackage{multicol}     % Gestion de la mise en page sur plusieurs colonnes

% Configuration du style des pieds de page
\pagestyle{fancy}
\fancyhf{} % Réinitialise les en-têtes et pieds de page
\renewcommand{\headrulewidth}{0pt}  % Pas de ligne sous l'en-tête
\renewcommand{\footrulewidth}{1pt}  % Ligne au-dessus du pied de page
\renewcommand{\footruleskip}{10pt}  % Espacement de la ligne de pied de page

%\fancyfoot[L]{\color{blue}\ding{45}\ \textbf{Prof: M. Diatta}} % Pied de page gauche
%\fancyfoot[R]{\color{blue}\ding{45}\ \textbf{2026}}        % Pied de page droit
%\cfoot{\bf\thepage\ / \pageref{LastPage}}                 % Numérotation centrale (X/Y)

% ==========================================
% 3. MATHÉMATIQUES ET SYSTÈMES
% ==========================================
\usepackage{amsmath}     % Environnements et commandes mathématiques de base
\usepackage{amssymb}     % Symboles mathématiques avancés
\usepackage{mathrsfs}    % Police calligraphique mathématique (\mathscr)

% Force l'affichage grand format (displaystyle) pour les formules intégrées au texte
\everymath{\displaystyle}

% ==========================================
% 4. TABLEAUX ET FIGURES
% ==========================================
\usepackage{array}       % Amélioration du traitement des tableaux
\usepackage{multirow}    % Fusion de cellules sur plusieurs lignes
\usepackage{float}        % Contrôle avancé du placement des figures ([H])
\usepackage{varwidth}    % Ajustement automatique de la largeur des boîtes

% ==========================================
% 5. LISTES ET PUCES PERSONNALISÉES
% ==========================================
\usepackage{enumitem}    % Personnalisation avancée des listes à puces et numérotées

% Styles TikZ pour puces circulaires et carrées
\newcommand\circitem[1]{%
  \tikz[baseline=(char.base)]{
    \node[circle, draw=gray, fill=red!55, minimum size=1.2em, inner sep=0] (char) {#1};
  }%
}
\newcommand\boxitem[1]{%
  \tikz[baseline=(char.base)]{
    \node[fill=cyan, minimum size=1.2em, inner sep=0] (char) {#1};
  }%
}

% Application aux niveaux 1 et 2 des listes numérotées
\setlist[enumerate,1]{label=\protect\circitem{\arabic*}}
\setlist[enumerate,2]{label=\protect\boxitem{\alph*}}

% ==========================================
% 6. GRAPHISMES, TIKZ ET TABLEAUX DE VARIATION
% ==========================================
\usepackage{tikz}        % Moteur de dessin vectoriel
\usetikzlibrary{patterns, trees} % Extensions TikZ :
                                 % - patterns : motifs de hachures et remplissages
                                 % - trees : réalisation d'arbres de dénombrement/probabilités

\usepackage{tkz-tab}     % Tableaux de variations et de signes mathématiques

% ==========================================
% 7. BOÎTES ET ENCADRÉS (TCOLORBOX)
% ==========================================
\usepackage[most]{tcolorbox} % Boîtes encadrées stylisées
\tcbuselibrary{skins, hooks} % Extensions graphiques (skins) et accroches d'actions (hooks)

% Environnement de boîte personnalisée pour les exemples
\newtcolorbox{exa}[2][]{
  enhanced, breakable,
  before skip=2mm, after skip=5mm,
  colback=yellow!20!white, colframe=black!20!blue, boxrule=0.5mm,
  attach boxed title to top left={xshift=0.6cm, yshift*=1mm-\tcboxedtitleheight},
  fonttitle=\bfseries,
  title={#2}, #1,
  boxed title style={
    frame code={
      \path[fill=tcbcolback!30!black]
        ([yshift=-1mm, xshift=-1mm]frame.north west)
        arc[start angle=0, end angle=180, radius=1mm]
        ([yshift=-1mm, xshift=1mm]frame.north east)
        arc[start angle=180, end angle=0, radius=1mm];
      \path[left color=tcbcolback!60!black, right color=tcbcolback!60!black, middle color=tcbcolback!80!black]
        ([xshift=-2mm]frame.north west) -- ([xshift=2mm]frame.north east)
        [rounded corners=1mm]-- ([xshift=1mm, yshift=-1mm]frame.north east)
        -- (frame.south east) -- (frame.south west)
        -- ([xshift=-1mm, yshift=-1mm]frame.north west)
        [sharp corners]-- cycle;
    },
    interior engine=empty,
  },
  interior style={top color=yellow!5}
}

% ==========================================
% 8. ARRIÈRE-PLAN ET FILIGRANE (WATERMARK)
% ==========================================
\usepackage{eso-pic}     % Ajout d'éléments en arrière-plan des pages

% Filigrane "PGB" centré sur chaque page
\AddToShipoutPicture{
  \AtTextCenter{%
    \makebox[0pt]{\rotatebox{50}{\textcolor[gray]{0.6}{\fontsize{1.5cm}{1.5cm}\selectfont maths-mes-maths}}}
  }
}

% ==========================================
% 9. COMPTEURS ET COMMANDES DE COURS
% ==========================================
% Compteur et commande pour la numérotation des exercices
\newcounter{exercice}
\renewcommand{\theexercice}{\arabic{exercice}}
\newcommand{\exo}{\refstepcounter{exercice}\textbf{Exercice \theexercice}\ }

% ==========================================
% 10. HYPERLIENS ET INTERACTIVITÉ
% ==========================================
% Chargé en tout dernier pour éviter les conflits de redéfinition de commandes
\usepackage[colorlinks=true, linkcolor=blue, urlcolor=blue, citecolor=blue]{hyperref}
\begin{document}
\renewcommand{\arraystretch}{1.5}
\renewcommand{\arrayrulewidth}{1.2pt}
\begin{tcolorbox}[
    colback=white,
    colframe=blue!60!black,
    boxrule=0.8pt,
    arc=3pt,
    left=3mm, right=3mm, top=2.5mm, bottom=2.5mm
]
    \begin{center}
        \textbf{\large Lycée de Diattacounda} \hfill \textbf{A.S. : 2026/2027} \\[1mm]
        \textbf{Mathématiques} $\bullet$ \textbf{Niveau : TS2} \hfill \textbf{Date : 13/11/2026} \\[2mm]
        \rule{0.8\linewidth}{0.4pt} \\[1.5mm]
        {\Large\bfseries\sffamily TD : Nombres Complexes}
    \end{center}
\end{tcolorbox}
\vspace{4mm}
%\begin{multicols}{2}
%\setlength{\columnseprule}{0.1mm} % La largeur de la ligne verticale entre les colonnes

\fbox{\textbf{\exo}}Démontrer que $(a, b, c, x, y \in \mathbb{R})$
\begin{enumerate}
\item \( (ax + by)^2 + (ay - bx)^2 = (a^2 + b^2)(x^2 + y^2) \)
\item \( (ax + by)^2 - (ay + bx)^2 = (a^2 - b^2)(x^2 - y^2) \)
\item \( (a + b + c)^3 - 3(a + b)(b + c)(c + a) = a^3 + b^3 + c^3 \)
\item \( (a - b)^3 + (b - c)^3 + (c - a)^3 = a^3 + b^3 + c^3 \)
\item \( (a^2 + b^2)(a'^2 + b'^2) = (aa' + bb')^2 + (ab' - ba')^2 \)
\item \( (a + b + c)^2 + (b - c)^2 + (c - a)^2 + (a - b)^2 = 3(a^2 + b^2 + c^2) \)
\item Soit $a, b, c$ trois réels non nul tels que $ab + bc + ca = 0$ \\
Calculer la somme $S=\frac{b+c}{a} + \frac{c+a}{b} + \frac{a+b}{c}$
\end{enumerate}
\fbox{\textbf{\exo}} 
\begin{enumerate}
\item Écrire sous la forme $2^m \times 3^n \times 5^p$ (avec $m, n, p$ des entiers relatifs) les réels suivants :
  \[
  A = \dfrac{(0{,}009)^{-3} \times (0{,}016)^2 \times 250}{(0{,}00075)^{-1} \times 810^3 \times 30}
  \quad ; \quad
  B = \dfrac{(-6)^4 \times 30^{-2} \times (-10)^{-3} \times 15^4}{(-25)^2 \times (36)^{-5} \times (-12)^3}
  \]
\item Écrire sous la forme $a^m b^n c^p$ (avec $m, n, p$ entiers relatifs)
  \[
  Q = \dfrac{(a^2 b)^{-3} \times (b c^3) \times (a^{-2} b^5)^3}{(b^2 c^2 a)^{-4} \times (a^{-1} b^6)^2}
  \quad ; \quad
  R = \dfrac{(a^{-2} c)^{-4} \times (-b^2 c)^5 \times (a^3 b c^{-1})^{-2}}{(-a^2 b^{-3} c)^3 \times (-b^4) \times (a^{-5} c)^2}
  \]
\item Calculer $AB$, $A^2 B^2$, $A^3 B^4$ où $A = -\left(\dfrac{1}{5}\right)^{-1} \times 3^2$
  et $B = \left(-\dfrac{1}{3}\right)^{-4} \times 5$
\end{enumerate}
\fbox{\textbf{\exo}} Simplifier les expressions suivantes :
\begin{enumerate}
\item \[
A = \dfrac{(-a)^7 \times (b^3 c^2)^4}{-b^3 c (-a)^4}
\quad ; \quad
B = \dfrac{91^{{-1}} \times (-39)^3 \times 25}{26^2 \times 45 \times (-21)^{-2} \times 72}
\quad ; \quad
C = \dfrac{(28 \times 12^{-2})^3 \times 105^{-3}}{7^{-2} \times (-60)^{-4} \times 63^4} \div \left(\dfrac{5}{3}\right)^5
\]
\item \[
D = \dfrac{\left[\left(-\dfrac{2}{3}\right)^2\right]^6 \times \left[\left(\dfrac{3}{5}\right)^{-2}\right]^3 \times \left[\left(\dfrac{5}{2}\right)^2\right]^{-3}}{\left(-\dfrac{4}{9}\right)^6}
\quad ; \quad
E = \dfrac{(a^2 b^3 c)^2 \times a^3 c}{(ab)^4 (c^2 a)^2 b c}
\]
\end{enumerate}

\fbox{\textbf{\exo}}
\begin{enumerate}
    \item Développer $(a + b + c)^2$.
    \item Montrer que si $a + b + c = 0$ alors $a^2 + b^2 + c^2 = -2 (ab + bc + ca)$.
    \item On suppose $a$, $b$ et $c$ sont non nuls. \\
    Montrer que $\frac{1}{a} + \frac{1}{b} + \frac{1}{c} = 0 \implies (a + b + c)^2 = a^2 + b^2 + c^2$.
\end{enumerate}
\fbox{\textbf{\exo}} Soit $a, b, c$ trois réels :
\begin{enumerate}
    \item Développer $(a + b + c)(ab + bc + ca)$ puis $(a + b + c)^3$
    \item Démontrer que si $a + b + c = 0$ alors $a^3 + b^3 + c^3 = 3abc$
    \item En déduire que, pour tous réel $x, y, z$ on a :
    $$(x - y)^3 + (y - z)^3 + (z - x)^3 = 3 (x - y) (y - z) (z - x)$$
\end{enumerate}

\fbox{\textbf{\exo}} Soient 4 entiers naturels consécutifs $n, n+1, n+2, n+3$, ($n > 0$)
\begin{enumerate}
    \item 
    \begin{enumerate}
        \item Démontrer que $(n+1) (n+2) = n (n+3) + 2$
        \item On pose $(n+1) (n+2) = a$. Exprimer en fonction de $a$ le produit $p = n (n+1) (n+2) (n+3)$
        \item En déduire que $p + 1$ est le carré d'un entier (on dit carré parfait)
    \end{enumerate}
    \item Déterminer $n$ sachant que $p = 5040$.
\end{enumerate}

\fbox{\textbf{\exo}}

\begin{enumerate}
    \item Calculer, pour tout entier $n$ supérieur ou égal à $2$ : $A = \left(1+\frac{1}{n-1}\right) \left(1-\frac{1}{n}\right)$
    \item Calculer $B = \left(1-\frac{1}{2^2}\right) \left(1-\frac{1}{3^2}\right) \left(1-\frac{1}{4^2}\right)$
    \item Calculer en fonction de $n$ : $X_n = \left(1-\frac{1}{2^2}\right) \left(1-\frac{1}{3^2}\right) \left(1-\frac{1}{4^2}\right) \dots \left(1-\frac{1}{n^2}\right)$
\end{enumerate}

\fbox{\textbf{\exo}}

Vérifier les identités suivantes :

\begin{enumerate}
    \item $(x+y+z)\left(\frac{1}{x} + \frac{1}{y} + \frac{1}{z}\right) = 1 + \frac{(x+y)(y+z)(z+x)}{xyz}$
    \item $\frac{x^2}{(x-t)(x-z)} + \frac{t^2}{(t-z)(t-x)} + \frac{z^2}{(z-x)(z-t)} = 1$
\end{enumerate}

\fbox{\textbf{\exo}}
\begin{enumerate}
    \item Calculer $A = (a^2+b^2+c^2)^2$
    \item Démontrer que : $(a+b+c=0) \implies [(a^2+b^2+c^2)^2 = 4(ab+bc+ca)^2]$
\end{enumerate}

\fbox{\textbf{\exo}}

Soit $x, y$ et $z$ trois réels tels que: $xyz = 1$. \\
Montrer que :
\[
\frac{x}{xy+x+1} + \frac{y}{yz+y+1} + \frac{z}{zx+z+1} = 1
\]

\fbox{\textbf{\exo}} Montrer que l'expression :
\[
E = \frac{4a^2-1}{(a-b)(a-c)} + \frac{4b^2-1}{(b-c)(b-a)} + \frac{4c^2-1}{(c-a)(c-b)} \text{ est égale à } 4.
\]

\fbox{\textbf{\exo}}
\begin{enumerate}
\item On donne : $a = \sqrt{x + \sqrt{x^2 - y^2}}$ et $b = \sqrt{x - \sqrt{x^2 - y^2}}$ ; avec $x > y \geq 0$.
\begin{enumerate}
    \item Calculer $a^2$, $b^2$ et $ab$. 
    \item En déduire que $(a+b)^2 = 2(x+y)$.
\end{enumerate}
\item Soient $x$ et $y$ deux nombres réels tels que $x > y \geq 0$. \\
Montrer que $\frac{\sqrt{x} - \sqrt{y}}{\sqrt{x} - \sqrt{y}} = \frac{\sqrt{x} + \sqrt{y}}{\sqrt{x} - \sqrt{y}}$
\item Soient $a$ et $b$ deux nombres réels strictement positifs.
\begin{enumerate}
    \item Développer l'expression $(\sqrt{a} - \sqrt{b})^2$.
    \item En déduire que $\sqrt{ab} \leq \frac{a+b}{2}$.
\end{enumerate}
\item  Simplifier $A = \sqrt{2} \times \sqrt{2+\sqrt{2}} \times \sqrt{2+\sqrt{2+\sqrt{2}}} \times \sqrt{2-\sqrt{2+\sqrt{2}}}$.
\end{enumerate}

\end{document}
